\subsection{Enrichment Scores}
\label{supp:enrichment-scores}

This section defines the DELT-Hit enrichment methods and the experimental
settings where each is most appropriate.

\textbf{Count Method.}
The count method averages raw sequencing counts across replicates within each
group and scores each compound by the difference between the condition
and control means, $\bar{c}_{\rm cond} - \bar{c}_{\rm ctrl}$. It is a
transparent count-based baseline that averages replicate counts and does not
assume a parametric count distribution. Its main weakness is sensitivity to
read-depth differences and raw-count scale, because the current default workflow
operates on raw counts rather than library-size-normalized frequencies. It is
therefore most informative when sequencing depths are comparable or as a simple
benchmark baseline.

\textbf{Z-Score Method (Faver et al.\ formulation).}
Let $D$ denote the total number of distinct disynthon (or higher-cycle) compound
identities in the library and $n$ the total number of sequencing reads in a
selection. For a given compound, let $c$ be its raw read count. The observed
sequence frequency is $p_{\rm obs} = c / n$ and the expected null frequency under
uniform library representation is $p_i = 1/D$. The depth-insensitive z-score
is\cite{faver_quantitativecomparisonenrichment_2019}

\begin{equation}
z_n = \frac{p_{\rm obs} - p_i}{\sqrt{p_i\,(1 - p_i)}}.
\label{eq:zscore}
\end{equation}

The closely related fold enrichment,
\begin{equation}
E = \frac{p_{\rm obs}}{p_i} = c \cdot \frac{D}{n},
\label{eq:fe}
\end{equation}
provides a scale-invariant measure of enrichment that is more interpretable for
hit-picking. Equations~\ref{eq:zscore} and \ref{eq:fe} summarize these two
related views. The recommended primary hit threshold is $E \geq 2.0$; a fixed
$z_n$ cutoff is not recommended because practical score ranges depend on
library diversity, sampling sparsity, and selection context.

\textbf{edgeR Method.}
Compound-level read counts per selection are modelled as negative-binomial
(NB) random variables. The NB model accommodates overdispersion relative to
Poisson counts, which arises in DEL experiments from variation in sequencing
depth, barcode-decoding efficiency, and chemical synthesis yield. Per-compound
dispersion is estimated by empirical Bayes shrinkage towards a trend fitted
across all compounds\cite{robinson_edgerbioconductorpackage_2010}. A likelihood-ratio test (LRT) is performed
for each compound to compare the NB-GLM fit under a two-group design
(condition vs.\ control) against a null model; Benjamini--Hochberg
false-discovery rate (FDR) correction is applied across all compounds.
Compounds are ranked by log$_2$ fold change (logFC); FDR is available as a
secondary filter.
